\section*{Chapter \ref{ch:b1:lab-techniques-1} --- Lab Techniques I: Safety, Measurement and Separation}

\begin{solution}{exo:b1:lab-techniques-1:1}
\emph{Danger}: H225 (highly flammable liquid, category 2) calls for it, and
the label carries the more severe word. H225 is a physical \omterm{def:b1:lab-techniques-1:hazard}{hazard}; H319
(serious eye irritation) and H336 (drowsiness or dizziness) are health
\omterm{def:b1:lab-techniques-1:hazard}{hazards}. Precautions: no flame or spark nearby, a closed bottle; goggles;
work in a ventilated place or a fume cupboard.
% ledger: ghs:acetone, hstat:H225, hstat:H319, hstat:H336
\end{solution}

\begin{solution}{exo:b1:lab-techniques-1:2}
(a) A \omterm{def:b1:lab-techniques-1:hazard}{hazard}, a property of the acid. (b) A \omterm{def:b1:lab-techniques-1:hazard}{risk}, small because the
exposure is small. (c) A \omterm{def:b1:lab-techniques-1:hazard}{hazard}. (d) A \omterm{def:b1:lab-techniques-1:hazard}{risk}: same \omterm{def:b1:lab-techniques-1:hazard}{hazard}, larger exposure
(vapour, splashes) when hot and open.
\end{solution}

\begin{solution}{exo:b1:lab-techniques-1:3}
$R_f = 2.1/6.0 = 0.35$ and $4.2/6.0 = 0.70$. The second spot has the same
$R_f$ as the reference: the sample may contain it (it is consistent), but
equal $R_f$ on one eluent do not prove identity; a second eluent, or a
co-spot of sample and reference that stays a single spot, strengthens the
conclusion. The spot at 0.35 is certainly another species.
\end{solution}

\begin{solution}{exo:b1:lab-techniques-1:4}
Half-width \qty{0.05}{mL}, rectangular: $u = 0.05/\sqrt3 = \qty{0.029}{mL}$.
For a difference of two independent readings,
$u = \sqrt{2} \times 0.029 = \qty{0.041}{mL}$.
\end{solution}

\begin{solution}{exo:b1:lab-techniques-1:5}
Mean \qty{0.25100}{g}; deviations $+2$, $-2$, $+5$, 0, $-5$ (in
\qty{e-4}{g}), sum of squares \qty{58e-8}{g^2}, $s = \sqrt{58 \times
10^{-8}/4} = \qty{3.8e-4}{g}$; $u = s/\sqrt5 = \qty{1.7e-4}{g}$;
$U = \qty{3.4e-4}{g}$: $m = (0.25100 \pm 0.00034)$~g, $k = 2$.
\end{solution}

\begin{solution}{exo:b1:lab-techniques-1:6}
One portion: $f = 1/(1 + 3 \times 30/50) = 0.357$, \qty{64.3}{\percent}
extracted. Two of \qty{15}{mL}: $f = 1.9^{-2} = 0.277$,
\qty{72.3}{\percent}. Three of \qty{10}{mL}: $f = 1.6^{-3} = 0.244$,
\qty{75.6}{\percent}.
\end{solution}

\begin{solution}{exo:b1:lab-techniques-1:7}
$z = 0.0012/\sqrt{0.0006^2 + 0.0002^2} = 0.0012/0.00063 = 1.9 \le 2$:
compatible, just. With \qty{0.0004}{mol/L}: $z = 0.0012/0.00045 = 2.7$,
not compatible: the more precise \omterm{def:b1:titration-methods:titration}{titration} reveals a real difference
(an error in the preparation, or in the \omterm{def:b1:titration-methods:titration}{titration}).
\end{solution}

\begin{solution}{exo:b1:lab-techniques-1:8}
Q: little soluble cold, very soluble hot. P dissolves the solid hardly
better hot than cold (nothing crystallises on cooling); R keeps too much
dissolved cold (large losses). With Q, \qty{5.0}{g} needs at least
$5.0/12 \times 100 = \qty{42}{mL}$ of boiling solvent, which keeps
$0.3 \times 0.417 = \qty{0.13}{g}$ dissolved cold: at most
\qty{4.9}{g}, \qty{97.5}{\percent}, can be recovered (before any loss on
the filter).
\end{solution}

\begin{solution}{exo:b1:lab-techniques-1:9}
Ethanol ($n_D = 1.3611$ at \qty{20}{\celsius}); water (1.333) and
cyclohexane (1.4266) are excluded. The index varies with temperature,
so a value means nothing without it. A value of 1.350, between water and
ethanol, suggests a mixture, such as ethanol containing water.
% ledger: nD:water, nD:ethanol, nD:cyclohexane
\end{solution}

\begin{solution}{exo:b1:lab-techniques-1:10}
With $t = KV/V_a$ and $n$ portions, $\ln f_n = -n\ln(1 + t/n)$; as
$n \to \infty$, $n \ln(1 + t/n) \to t$ (since $\ln(1 + \varepsilon) \sim
\varepsilon$), and $f_n$ decreases towards $\mathrm{e}^{-t}$ (\cref{prop:b1:lab-techniques-1:multiple-extractions}).
Here $t = 3 \times 30/50 = 1.8$, $\mathrm{e}^{-1.8} = 0.165$: at most
\qty{83.5}{\percent}. Three portions already give \qty{75.6}{\percent},
\qty{91}{\percent} of the limit; reaching \qty{99}{\percent} of it takes 32
portions of less than \qty{1}{mL}: beyond two or three portions, the gain is
not worth the work.
\end{solution}

\begin{solution}{exo:b1:lab-techniques-1:11}
With $t = x - x_0$: $\int_{-a}^{a} (a - |t|)/a^2 \,\mathrm{d}t =
2 \int_0^a (a - t)/a^2 \,\mathrm{d}t = 2 (a^2/2)/a^2 = 1$. The mean is
$x_0$ by symmetry; the variance is $2 \int_0^a t^2 (a - t)/a^2
\,\mathrm{d}t = \frac{2}{a^2}\left(\frac{a^4}{3} - \frac{a^4}{4}\right) =
\frac{a^2}{6}$, so $u = a/\sqrt6 = 0.41\,a$, against $a/\sqrt3 = 0.58\,a$
for the rectangular case: knowing that the middle is more likely reduces
the uncertainty.
\end{solution}

\begin{solution}{exo:b1:lab-techniques-1:12}
$c = m/(MV) = 0.25100/(204.2 \times 0.10000) = \qty{0.012292}{mol/L}$.
Relative uncertainties: mass $1.7 \times 10^{-4}/0.2510 = 6.8 \times
10^{-4}$, volume $0.06/100.00 = 6.0 \times 10^{-4}$; combined
$\sqrt{(6.8^2 + 6.0^2)} \times 10^{-4} = 9.1 \times 10^{-4}$, that is
\qty{0.091}{\percent}; $u(c) = \qty{1.1e-5}{mol/L}$, $U = \qty{2.2e-5}{mol/L}$:
$c = (0.012292 \pm 0.000022)$~mol/L, $k = 2$. The weighing dominates,
slightly; the two sources are of the same size.
\end{solution}

\begin{solution}{pb:b1:lab-techniques-1:1}
\textbf{1.} GHS02 flammable, GHS05 corrosive, GHS06 acute toxicity, GHS07
harmful or irritant.
\textbf{2.} Salicylic acid: \emph{Danger}, because of H318 (serious eye
damage, category 1). Aspirin: \emph{Warning} (H302 only).
\textbf{3.} The \omterm{def:b1:lab-techniques-1:hazard}{hazard} is fixed, the exposure is not: a small volume,
measured and used in a fume cupboard, with gloves, goggles and a lab coat,
the bottle closed at once, no flame nearby.
\textbf{4.} H314 (severe skin burns and eye damage): goggles and gloves.
H332 (harmful if inhaled) and H226 (flammable liquid and vapour): the fume
cupboard and the absence of flames.
\textbf{5.} Rinse cautiously with water for several minutes, removing
contact lenses if possible, and keep rinsing (P305+P351+P338); then seek
medical advice.
\textbf{6.} Into the aqueous waste, after neutralisation, not down the
sink.
\textbf{7.} \ce{C7H6O3 + C4H6O3 -> C9H8O4 + C2H4O2}: C $7 + 4 = 9 + 2$,
H $6 + 6 = 8 + 4$, O $3 + 3 = 4 + 2$.
\textbf{8.} $2.00/138.0 = \qty{1.449e-2}{mol}$; at most
$\num{1.449e-2} \times 180.0 = \qty{2.61}{g}$ of aspirin.
\textbf{9.} The crude solid contains unreacted salicylic acid, ethanoic
acid and water: its mass is not that of aspirin, and its purity is
unknown.
\textbf{10.} What stays dissolved after cooling is lost: a \omterm{def:b1:precipitation:solubility}{solubility} small
cold and large hot means the product dissolves when hot and comes back
almost entirely when cold.
\textbf{11.} Every extra millilitre keeps its share of product dissolved
when cold.
\textbf{12.} Slow growth gives larger, more regular \omterm{def:b1:crystals-metals:crystal}{crystals} that trap
less mother liquor and impurities; the ice lowers the \omterm{def:b1:precipitation:solubility}{solubility}, hence
the loss.
\textbf{13.} $\qty{0.045}{L} \times \qty{3.3}{g/L} = \qty{0.15}{g}$.
\textbf{14.} $1.95/2.61 = \qty{75}{\percent}$ (\qty{74.7}{\percent}).
\textbf{15.} $807/6 = \qty{134.5}{\celsius}$.
\textbf{16.} Deviations $-0.5$, $0.5$, $-1.5$, $0.5$, $-0.5$, $1.5$; sum of
squares 5.5; $s = \sqrt{5.5/5} = \qty{1.05}{\celsius}$.
\textbf{17.} $u_A = 1.05/\sqrt6 = \qty{0.43}{\celsius}$.
\textbf{18.} Half-width \qty{0.5}{\celsius}: $0.5/\sqrt3 = \qty{0.29}{\celsius}$.
\textbf{19.} $1/\sqrt3 = \qty{0.58}{\celsius}$;
$u = \sqrt{0.428^2 + 0.289^2 + 0.577^2} = \sqrt{0.600} = \qty{0.77}{\celsius}$.
\textbf{20.} $U = 2 \times 0.775 = \qty{1.5}{\celsius}$:
$T = (134.5 \pm 1.5)$~\unit{\celsius}, $k = 2$.
\textbf{21.} It melts low and over a range of \qty{8}{\celsius}: it is
impure (salicylic acid, ethanoic acid, water).
\textbf{22.} Given to the unit, the value lies within $\pm\qty{0.5}{\celsius}$:
$u_{\mathrm{ref}} = 0.5/\sqrt3 = \qty{0.29}{\celsius}$.
\textbf{23.} $2.3/6.0 = 0.38$ and $3.3/6.0 = 0.55$. The crude product
contained salicylic acid (same $R_f$ as the reference) and aspirin; after
\omterm{def:b1:lab-techniques-1:recrystallisation}{recrystallisation} only the aspirin spot remains: the salicylic acid has
been removed, at least below what the plate detects.
\textbf{24.} Aspirin decomposes as it melts; heated slowly, it has time to
decompose, and the mixture formed melts lower. The tabulated value refers
to rapid heating, as on the bench.
\textbf{25.} $z = |134.5 - 135|/\sqrt{0.600 + 0.083} = 0.5/0.827 =$
\textbf{0.60}: well below 2, the measured and tabulated melting points are
compatible.
\end{solution}
